% titel: Blatt 8, A6: Natürlicher Logarithmus
% Teil II: der durchgerechnete Weg.

\subsection*{Gegeben und gesucht}
$L\colon(0,\infty)\to\mathbb R$ differenzierbar mit
$L(1)=0$, $L'(1)=1$ und $L(xy)=L(x)+L(y)$ für alle $x,y>0$. Zeige:
$L=\ln$.

\subsection*{Differentialgleichung aus der Funktionalgleichung}
Differenziere $L(xy)=L(x)+L(y)$ nach $y$; dabei ist $x$ fest, links wendet
man die Kettenregel an ($\frac d{dy}(xy)=x$):
\begin{equation*}
  L'(xy)\cdot x=L'(y).
\end{equation*}
Setze nun $y=1$:
\begin{equation*}
  x\,L'(x)=L'(1)=1\;\Longrightarrow\;
  L'(x)=\frac1x\qquad\forall x>0 .
\end{equation*}

\subsection*{Integration}
Eine Stammfunktion von $\frac1x$ ist $\ln x$, und $L(1)=0=\ln1$. Aus
$L'=\ln'=\frac1x$ und $L(1)=\ln1=0$ folgt
\begin{equation*}
  L(x)=\ln x\qquad\forall x>0 .
\end{equation*}

\begin{ergebnis}
$L(x)=\ln x$ für alle $x>0$. Der drei Eigenschaften genügen nur genau die
Logarithmusfunktion.
\end{ergebnis}

\begin{kontrolle}
\textbf{Probe:} $\ln(xy)=\ln x+\ln y$, $\ln 1=0$, $\ln'1=\frac11=1$ -- die
drei axiomatischen Eigenschaften sind für $\ln$ erfüllt, und die Herleitung
zeigt, dass es keine anderen Funktionen gibt. $\ln'(x)=\frac1x$ ist die
einzige Lösung mit $\ln 1=0$.

\textbf{Grenzfall:} Ohne $L'(1)=1$ wäre jede Konstante $c$ möglich und
$L(x)=c\ln x$; die Bedingung $L'(1)=1$ legt $c=1$ fest.
\end{kontrolle}

\begin{merksatz}
Die Additionseigenschaft $L(xy)=L(x)+L(y)$ ist die definierende Identität
des Logarithmus. Ableiten nach $y$ gibt $xL'(xy)=L'(y)$; mit $y=1$
entsteht $L'(x)=1/x$; mit $L(1)=0$ folgt $L(x)=\ln x$. Die Bedingung
$L'(1)=1$ fixiert die Normierung.
\end{merksatz}
